% Options for packages loaded elsewhere
\PassOptionsToPackage{unicode}{hyperref}
\PassOptionsToPackage{hyphens}{url}
\PassOptionsToPackage{dvipsnames,svgnames,x11names}{xcolor}
\documentclass[
  11pt,
]{article}
\usepackage{xcolor}
\usepackage[margin=27mm]{geometry}
\usepackage{amsmath,amssymb}
\setcounter{secnumdepth}{2}
\usepackage{iftex}
\ifPDFTeX
  \usepackage[T1]{fontenc}
  \usepackage[utf8]{inputenc}
  \usepackage{textcomp} % provide euro and other symbols
\else % if luatex or xetex
  \usepackage{unicode-math} % this also loads fontspec
  \defaultfontfeatures{Scale=MatchLowercase}
  \defaultfontfeatures[\rmfamily]{Ligatures=TeX,Scale=1}
\fi
\usepackage{lmodern}
\ifPDFTeX\else
  % xetex/luatex font selection
\fi
% Use upquote if available, for straight quotes in verbatim environments
\IfFileExists{upquote.sty}{\usepackage{upquote}}{}
\IfFileExists{microtype.sty}{% use microtype if available
  \usepackage[]{microtype}
  \UseMicrotypeSet[protrusion]{basicmath} % disable protrusion for tt fonts
}{}
\makeatletter
\@ifundefined{KOMAClassName}{% if non-KOMA class
  \IfFileExists{parskip.sty}{%
    \usepackage{parskip}
  }{% else
    \setlength{\parindent}{0pt}
    \setlength{\parskip}{6pt plus 2pt minus 1pt}}
}{% if KOMA class
  \KOMAoptions{parskip=half}}
\makeatother
\usepackage{longtable,booktabs,array}
\newcounter{none} % for unnumbered tables
\usepackage{calc} % for calculating minipage widths
% Correct order of tables after \paragraph or \subparagraph
\usepackage{etoolbox}
\makeatletter
\patchcmd\longtable{\par}{\if@noskipsec\mbox{}\fi\par}{}{}
\makeatother
% Allow footnotes in longtable head/foot
\IfFileExists{footnotehyper.sty}{\usepackage{footnotehyper}}{\usepackage{footnote}}
\makesavenoteenv{longtable}
\setlength{\emergencystretch}{3em} % prevent overfull lines
\providecommand{\tightlist}{%
  \setlength{\itemsep}{0pt}\setlength{\parskip}{0pt}}
\usepackage{bookmark}
\IfFileExists{xurl.sty}{\usepackage{xurl}}{} % add URL line breaks if available
\urlstyle{same}
\hypersetup{
  colorlinks=true,
  linkcolor={Maroon},
  filecolor={Maroon},
  citecolor={Blue},
  urlcolor={Blue},
  pdfcreator={LaTeX via pandoc}}

\title{Complete monotonicity of Tsallis entropy under heat flow:\\a sharp order classification}
\author{Anonymous}
\date{6 November 2026\\Version 0.1.0-candidate\\\small DOI: \href{https://doi.org/10.5281/zenodo.24001602}{10.5281/zenodo.24001602}}

\begin{document}

\maketitle

AI-assisted, unrefereed mathematical candidate. Original prose is
dedicated to the public domain under CC0 1.0; accompanying code is
supplied under its stated licence.

\begin{abstract}

For a probability density evolving under the linear heat equation,
consider the alternating signs of all derivatives of its Tsallis
entropy. We prove that the quadratic order \(\alpha=2\) is the unique
positive order for which these inequalities hold universally. For each
fixed \(\alpha>0\) with \(\alpha\notin\{1,2\}\), we give an explicit
family of smooth, strictly positive, strictly log-concave densities
\(f_{\alpha,m}\) with Gaussian tails, indexed by the derivative order
\(m\), for which the signed \(m\)th entropy derivative is negative for
every sufficiently large \(m\). The input density varies with \(m\).
Previously established Shannon counterexamples supply the case
\(\alpha=1\). The proof extends the two-frequency spectral construction
and Gaussian localisation of Zou, Fan, Gao and Wang: an order-dependent
choice of the second Fourier amplitude makes the unique dominant
spectral coefficient equal to \(-\alpha|\alpha-2|/8\). A separate
example at \(\alpha=3/2\) specifies a single real-line density and
certifies failure at derivative order \(13\) using an exact rational
periodic calculation and a quantitative analytic transfer. The result
classifies the order parameter pointwise; it does not claim the first
refutation of the Shannon conjecture or a new localisation mechanism.

\end{abstract}

\section{Question, conventions and
scope}\label{question-conventions-and-scope}

Adding Gaussian noise smooths a probability density. Entropy
inequalities ask which features of that smoothing can be described by
universal signs. An increasing entropy curve is one such property;
alternating signs at every derivative order are a much stronger
requirement. This paper determines precisely which positive Tsallis
orders can satisfy that whole hierarchy for every input density.

Let

\[
P_t=e^{t\Delta/2},\qquad t\geq0,
\]

be the heat semigroup on \(\mathbb R^d\). Thus \(P_tf\) is the density
obtained by adding centred Gaussian noise with covariance \(tI_d\). For
\(\alpha>0\), \(\alpha\ne1\), define

\[
\tag{1}
 J_{\alpha,f}(t)=\int_{\mathbb R^d}(P_tf)^\alpha\,\mathrm dx,
 \qquad
 S_{\alpha,f}(t)=\frac{1-J_{\alpha,f}(t)}{\alpha-1}.
\]

At \(\alpha=1\), use the continuous extension

\[
S_{1,f}(t)=-\int_{\mathbb R^d}(P_tf)\log(P_tf)\,\mathrm dx.
\]

The property at issue is

\[
\tag{2}
 (-1)^{k-1} S_{\alpha,f}^{(k)}(t)\geq0
 \quad\text{for every }k\geq1\text{ and }t>0.
\]

The first inequality says that entropy increases under smoothing; the
second says that its graph is concave. Higher inequalities impose
additional restrictions that do not follow from these two. Failure of
the full hierarchy therefore does not by itself refute an established
low-order sign theorem. Equivalently, \(S_{\alpha,f}'\) is completely
monotone. Entropy itself need not be nonnegative, so we state the
derivative inequalities directly. A derivative assertion is understood
where the relevant entropy is finite and the derivative exists; every
density constructed here has Gaussian tails and an entropy profile
analytic in a neighbourhood of the time being used.

Wu, Yu and Guo~\href{https://arxiv.org/abs/2312.01819v4}{{[}1{]}}
conjectured (2) for every \(\alpha\in(1,2)\), under this heat
convention, and proved the quadratic case. Their conjecture implies the
corresponding Shannon conjecture as \(\alpha\to1\). Gu and
Sellke~\href{https://arxiv.org/abs/2605.11656v1}{{[}2{]}} subsequently
disproved the Shannon conjecture by a certified finite atomic example
and also proved the existence of log-concave counterexamples. Zou, Fan,
Gao and Wang~\href{https://arxiv.org/abs/2608.30275v1}{{[}3{]}} then
gave an explicit smooth strictly log-concave Shannon family, using two
interacting Fourier modes localised by a Gaussian envelope.

Earlier results also delimit the scope of the question. Wu, Yu and
Guo~\href{https://arxiv.org/abs/2312.01819v4}{{[}1, Example 2{]}}
reported failures for the pairs
\((k,\alpha)=(3,33),(4,8),(5,5),(6,4),(9,3)\) using a symmetric
two-point input. Wu and
Yu~\href{https://arxiv.org/abs/2504.01837v2}{{[}4, Section III-B and
Appendix B{]}} established related low-order signs and conditional
implications, and gave the quadratic identity in arbitrary dimension.
Sun~\href{https://arxiv.org/abs/2601.12944v3}{{[}5{]}} proved concavity
of Tsallis entropy for \(\alpha\in(0,3]\) in dimension one and for
\(\alpha\in[1,3]\) in higher dimensions. These lower-order conclusions
are consistent with the higher-derivative failures proved here.

These results already preclude the entire interval conjecture from being
true: strict failure for Shannon entropy and continuity in \(\alpha\)
produce failures at some orders close to \(1\). The question treated
here is sharper: \emph{which individual positive orders can satisfy (2)
universally?} Our answer is that only \(\alpha=2\) survives, even if the
universal quantifier is restricted to smooth strictly positive strictly
log-concave densities with Gaussian tails.

The mechanism of the proof is adapted directly
from~\href{https://arxiv.org/abs/2608.30275v1}{{[}3{]}}, and we credit
that source throughout. The additional steps are the explicit tuning of
the second Fourier amplitude for an arbitrary fixed Tsallis order, and a
transfer estimate for the power moment, whose Gaussian scaling factor is
absent from the Shannon periodic term. The certified finite example is a
separate concrete witness; it is not a claim about the least possible
derivative order.

\textbf{Theorem 1} (Sharp universal classification). Fix a dimension
\(d\geq1\). The only \(\alpha>0\) for which (2) holds for every smooth
strictly positive strictly log-concave probability density with Gaussian
tails on \(\mathbb R^d\) is \(\alpha=2\).

For \(\alpha=2\), the inequalities hold, more generally, for every
probability density and every positive time. For each \(\alpha\ne2\), a
density in the stated smaller class violates a strict inequality at a
positive time.

The construction for \(\alpha\notin\{1,2\}\) is explicit. Fix such an
\(\alpha\), and set

\[
\tag{3}
 P_\alpha=(\alpha-3)(\alpha-4)(\alpha-5),\qquad
 B_\alpha=\operatorname{sgn}(2-\alpha)-\frac{P_\alpha}{288}.
\]

For integers \(m\geq1\), put

\[
\tag{4}
 a_m=2^{-m/4},\qquad b_m=B_\alpha2^{-m}=B_\alpha a_m^4,
 \qquad \kappa_m=8a_m,
\]

and, whenever the periodic factor is positive, define

\[
\tag{5}
 f_{\alpha,m}(x)=Z_{\alpha,m}^{-1}e^{-\kappa_m x^2/2}
       (1+a_m\cos x+b_m\cos2x),\qquad x\in\mathbb R.
\]

The exact normalising constant is

\[
\tag{6}
 Z_{\alpha,m}=\sqrt{\frac{2\pi}{\kappa_m}}\,\Xi_{\alpha,m},
 \qquad
 \Xi_{\alpha,m}=1+a_me^{-1/(2\kappa_m)}+b_me^{-2/\kappa_m}.
\]

The index \(m\) specifies both the density and the derivative order.

\textbf{Theorem 2} (Explicit log-concave family). Fix \(\alpha>0\),
\(\alpha\notin\{1,2\}\), and let (3)--(6) hold. For every sufficiently
large \(m\), \(f_{\alpha,m}\) is a smooth strictly positive strictly
log-concave probability density with Gaussian tails. Define

\[
\tag{7}
 \beta=\frac{3}{2\sqrt2}>1,\qquad
 C_{\alpha,m}=\alpha^{-1/2}
       \left(\frac{\kappa_m}{2\pi}\right)^{(\alpha-1)/2}
       \Xi_{\alpha,m}^{-\alpha}>0.
\]

Then

\[
\tag{8}
 \frac{(-1)^{m-1}S_{\alpha,f_{\alpha,m}}^{(m)}(0)}{C_{\alpha,m}}
 =-\frac{\alpha|\alpha-2|}{8}\,\beta^m+o_\alpha(\beta^m)
 \qquad(m\to\infty).
\]

In particular, for each sufficiently large \(m\), the signed derivative
is negative at zero and at every sufficiently small positive time.

All asymptotic statements in this theorem hold for a \emph{fixed}
\(\alpha\). No uniformity as \(\alpha\to1\), \(\alpha\to2\),
\(\alpha\to0\) or \(\alpha\to\infty\) is asserted.

\clearpage
\subsection{Witness map and proof
roadmap}\label{witness-map-and-proof-roadmap}

{\def\LTcaptype{none} % do not increment counter
\begin{longtable}[]{@{}
  >{\raggedright\arraybackslash}p{(\linewidth - 4\tabcolsep) * \real{0.3333}}
  >{\raggedright\arraybackslash}p{(\linewidth - 4\tabcolsep) * \real{0.3333}}
  >{\raggedright\arraybackslash}p{(\linewidth - 4\tabcolsep) * \real{0.3333}}@{}}
\toprule\noalign{}
\begin{minipage}[b]{\linewidth}\raggedright
Construction
\end{minipage} & \begin{minipage}[b]{\linewidth}\raggedright
Entropy order and derivative
\end{minipage} & \begin{minipage}[b]{\linewidth}\raggedright
Shape and scope
\end{minipage} \\
\midrule\noalign{}
\endhead
\bottomrule\noalign{}
\endlastfoot
Theorem 2 & Each fixed \(\alpha\notin\{1,2\}\); all sufficiently large
\(m\) & A different smooth strictly log-concave density for each \(m\);
threshold not quantified \\
Inherited Shannon theorem & \(\alpha=1\); all sufficiently large \(m\) &
The cited log-concave family, with density depending on \(m\) \\
Theorem 7 & \(\alpha=3/2\), derivative \(13\), time \(1/50\) & One
specified smooth positive density; no log-concavity or least-order
claim \\
\end{longtable}
}

The periodic calculation first arranges a negative coefficient at one
spectral level. Repeated differentiation amplifies that level faster
than all others. Gaussian localisation then turns the periodic
perturbation into an integrable probability density, while controlling
the errors even after the large number of derivatives. The quadratic
case is different: its power integral is a nonnegative Fourier-weighted
mixture of decaying exponentials. The finite witness later replaces an
unspecified asymptotic threshold by exact rational bounds for one order
and time.

\section{The periodic calculation}\label{the-periodic-calculation}

For a periodic integrable function, write

\[
\left\langle F\right\rangle_{\mathbb T}=\frac1{2\pi}\int_0^{2\pi}F(x)\,\mathrm dx.
\]

Define

\[
\tag{9}
 r_m(s,x)=1+a_me^{-s/2}\cos x+b_me^{-2s}\cos2x,
 \qquad
 D_{\alpha,m}(s)=\frac{\left\langle r_m(s,\cdot)^\alpha\right\rangle_{\mathbb T}-1}{\alpha-1}.
\]

This is the periodic part of the negative Tsallis entropy. The heat
exponents and the choice of amplitudes in (4) follow the spectral
construction in~\href{https://arxiv.org/abs/2608.30275v1}{{[}3{]}}.

\textbf{Lemma 3} (Periodic asymptotic). For every fixed \(\alpha>0\),
\(\alpha\notin\{1,2\}\),

\[
\tag{10}
 (-1)^mD_{\alpha,m}^{(m)}(0)
 =-\frac{\alpha|\alpha-2|}{8}\,\beta^m+O_\alpha(1).
\]

Moreover, the absolute sum of the differentiated monomial expansion is
\(O_\alpha(\beta^m)\).

\emph{Proof.} For large \(m\), \(a_m+|b_m|<1/2\). Expanding the analytic
function \((1+u)^\alpha\) around \(u=0\), using that the linear term has
zero circle mean, gives

\[
\tag{11}
 D_{\alpha,m}(s)=
 \sum_{\substack{p,q\geq0\\p+q\geq2}}
  \Gamma_{p,q}(\alpha)\,a_m^p b_m^q e^{-(p+4q)s/2},
\]

where

\[
\tag{12}
 \Gamma_{p,q}(\alpha)=
 \frac{\binom\alpha{p+q}}{\alpha-1}
 \binom{p+q}{p}\left\langle \cos^p x\cos^q2x\right\rangle_{\mathbb T}.
\]

For each fixed \(m\), this series can be differentiated \(m\) times near
zero: at total degree \(\ell=p+q\), the time factor contributes at most
\((2\ell)^m\), whereas the remaining absolute sum is bounded by a fixed
geometric factor times a binomial-series coefficient. Taking \(m\) large
enough that \(a_m+|b_m|\) is smaller than a fixed analytic radius yields
an absolutely convergent differentiated series.

The substitution \(x\mapsto\pi-x\) shows that the trigonometric average
in (12) vanishes for odd \(p\). Write \(p=2h\) and set \(j=h+2q\). Then

\[
\frac{p+4q}{2}=j,\qquad
 a_m^{2h}b_m^q=B_\alpha^q2^{-mj/2}.
\]

Consequently,

\[
\tag{13}
 (-1)^m D_{\alpha,m}^{(m)}(0)
 =\sum_{j\geq1} A_j(\alpha)\beta_j^m,
 \qquad
 \beta_j=j2^{-j/2},
\]

where

\[
A_j(\alpha)=
 \sum_{\substack{h,q\geq0\\h+2q=j\\2h+q\geq2}}
    \Gamma_{2h,q}(\alpha)B_\alpha^q.
\]

The ratio \(\beta_{j+1}/\beta_j=(j+1)/(j\sqrt2)\) shows that \(j=3\) is
the unique maximum. Specifically, \(\beta_2=\beta_4=1\),
\(\beta_3=\beta\), and \(\beta_j<1\) for every \(j\notin\{2,3,4\}\).

There are exactly two pairs at level \(j=3\): \((p,q)=(2,1)\) and
\((6,0)\). Since

\[
\left\langle \cos^2x\cos 2x\right\rangle_{\mathbb T}=\frac14,\qquad
 \left\langle \cos^6x\right\rangle_{\mathbb T}=\frac5{16},
\]

their sum is

\[
\begin{aligned}
 A_3(\alpha)
 &=\frac34\frac{\binom\alpha3}{\alpha-1}B_\alpha
   +\frac5{16}\frac{\binom\alpha6}{\alpha-1}\\
 &=\frac{\alpha(\alpha-2)}8 B_\alpha
   +\frac{\alpha(\alpha-2)(\alpha-3)(\alpha-4)(\alpha-5)}{2304}
   =-\frac{\alpha|\alpha-2|}{8}.
\end{aligned}\tag{14}
\]

This is the only order-dependent tuning needed to obtain a negative
leading coefficient.

For completeness, we control the infinite sum. Analyticity of

\[
\frac{(1+u)^\alpha-1-\alpha u}{\alpha-1}
\]

on \(|u|<1\) implies, by a Cauchy estimate on \(|u|=1/2\), that

\[
\left|\frac{\binom\alpha\ell}{\alpha-1}\right|
 \leq K_\alpha2^\ell,\qquad \ell\geq2.
\]

At level \(j\) there are at most \(j+1\) pairs and \(\ell=2j-3q\leq2j\).
Therefore the sum \(\mathcal A_j\) of the absolute monomial coefficients
at that level satisfies

\[
\tag{15}
 |A_j|\leq\mathcal A_j\leq
 K_\alpha(j+1)R_\alpha^j,
 \qquad R_\alpha=16\max\{1,|B_\alpha|\}^2.
\]

For \(j\geq16\), one has \(j\leq2^{j/4}\), hence
\(\beta_j\leq2^{-j/4}\). Choose an integer \(m_0\) with
\(R_\alpha2^{-m_0/4}<1\). Then

\[
\sum_{j\geq16}\mathcal A_j\beta_j^m
 \leq K_\alpha\sum_{j\geq16}(j+1)
       (R_\alpha2^{-m_0/4})^j<\infty,
 \qquad m\geq m_0,
\]

uniformly in \(m\). The finitely many levels \(j\leq15\), \(j\ne3\),
contribute \(O_\alpha(1)\). Equation (14) proves (10); retaining
absolute values proves the last assertion.~\(\square\)

\section{Gaussian localisation for power
moments}\label{gaussian-localisation-for-power-moments}

We now adapt the exact localisation method
in~\href{https://arxiv.org/abs/2608.30275v1}{{[}3{]}} to
\(J_{\alpha,f}\). Throughout this section, \(\alpha>0\),
\(\alpha\notin\{1,2\}\), is fixed, and the subscript \(m\) is sometimes
suppressed.

\subsection{Positivity and strict
log-concavity}\label{positivity-and-strict-log-concavity}

For sufficiently large \(m\), \(a_m+|b_m|\leq1/2\) and
\(|b_m|\leq a_m/8\). A simple sufficient condition is \(a_m\leq1/4\)
together with \(|B_\alpha|a_m^3\leq1/8\). If

\[
\rho_m(x)=1+a_m\cos x+b_m\cos2x,
\]

then \(\rho_m\geq1/2\) and \(|\rho_m''|\leq a_m+4|b_m|\leq3a_m/2\).
Therefore

\[
(\log\rho_m)''=\frac{\rho_m''}{\rho_m}
   -\frac{(\rho_m')^2}{\rho_m^2}\leq3a_m,
\]

and

\[
\tag{16}
 (\log f_{\alpha,m})''=-\kappa_m+(\log\rho_m)''
 \leq-5a_m<0.
\]

The remaining regularity and tail properties follow immediately from (5)
and the bounds on \(\rho_m\).

\subsection{Exact heat evolution}\label{exact-heat-evolution}

For \(\kappa>0\) and \(n\in\mathbb Z\), completing the square in the
Gaussian convolution gives the identity
from~\href{https://arxiv.org/abs/2608.30275v1}{{[}3, Lemma 3{]}}:

\[
\tag{17}
 P_t\big(e^{-\kappa x^2/2}e^{inx}\big)
 =(1+\kappa t)^{-1/2}
  e^{-\kappa x^2/(2(1+\kappa t))}
  e^{-n^2t/(2(1+\kappa t))}e^{inx/(1+\kappa t)}.
\]

Indeed, the exponent in the convolution integrand is

\[
-\frac{1+\kappa t}{2t}
   \left(y-\frac{x+int}{1+\kappa t}\right)^2
 -\frac{\kappa x^2}{2(1+\kappa t)}
 -\frac{n^2t}{2(1+\kappa t)}+\frac{inx}{1+\kappa t},
\]

and Gaussian integration proves (17) for \(t>0\), with continuity at
zero.

Set

\[
\tag{18}
 d_m(t)=1+\kappa_mt,\qquad
 \tau_m(t)=\frac{t}{d_m(t)},\qquad
 \lambda_m(t)=\alpha\kappa_m d_m(t),\qquad
 \nu=\frac{1-\alpha}{2}.
\]

Applying (17) to the three modes in (5) yields

\[
\tag{19}
 P_tf_{\alpha,m}(x)=Z_{\alpha,m}^{-1}d_m(t)^{-1/2}
 e^{-\kappa_mx^2/(2d_m(t))}
 r_m\!\left(\tau_m(t),\frac{x}{d_m(t)}\right).
\]

For real \(\lambda>0\), let \(\langle\cdot\rangle_\lambda\) denote
expectation under the centred Gaussian law of precision \(\lambda\).
Changing variables \(x=d_m(t)y\) in the power integral gives the scalar
identity

\[
Z_{\alpha,m}^{-\alpha}d_m(t)^{1-\alpha/2}
 \sqrt{\frac{2\pi}{\alpha\kappa_m d_m(t)}}
 =C_{\alpha,m}d_m(t)^{(1-\alpha)/2}.
\]

Consequently,

\[
\tag{20}
 J_{\alpha,f_{\alpha,m}}(t)
 =C_{\alpha,m}d_m(t)^\nu
  \left\langle r_m(\tau_m(t),\cdot)^\alpha
  \right\rangle_{\lambda_m(t)}.
\]

The constant \(C_{\alpha,m}\) is exactly (7); no asymptotic
normalisation is being substituted here.

\textbf{Lemma 4} (Power-moment transfer). As \(m\to\infty\) with
\(\alpha\) fixed,

\[
\tag{21}
 \frac{(-1)^mJ_{\alpha,f_{\alpha,m}}^{(m)}(0)}
 {C_{\alpha,m}(\alpha-1)}
 =(-1)^mD_{\alpha,m}^{(m)}(0)+o_\alpha(\beta^m).
\]

The profile in (20) has a holomorphic extension to a neighbourhood of
zero for all sufficiently large \(m\).

\emph{Proof.} We separate Fourier filtering, the scalar Gaussian factor,
and the time change. A broad Gaussian average suppresses nonconstant
Fourier modes. The power moment also introduces a scalar Gaussian
factor, absent from the Shannon periodic term, and a small
reparametrisation of time. Each contribution must remain smaller than
the negative \(\beta^m\) scale after taking \(m\) derivatives.

\paragraph{Fourier filtering.}\label{fourier-filtering.}

Use the complex time disk \(|t|\leq1/2\) and the spatial strip
\(|\operatorname{Im}y|\leq1\). For large \(m\), \(|\kappa_mt|\leq1/4\),
\(|\tau_m(t)|\leq1\), and

\[
\left|a_me^{-\tau_m(t)/2}\cos y+
       b_me^{-2\tau_m(t)}\cos2y\right|<\frac12
\]

on this product domain. Consequently

\[
G_m(t,y)=r_m(\tau_m(t),y)^\alpha
\]

is defined there by the principal power and is holomorphic and uniformly
bounded, with a bound depending only on \(\alpha\). Its Fourier
coefficients satisfy

\[
\tag{22}
 |\widehat G_m(n,t)|\leq K_\alpha e^{-|n|}.
\]

For complex \(\lambda\) with positive real part, the Gaussian bracket
has a holomorphic extension, using the principal square root, and

\[
\left\langle e^{iny}\right\rangle_\lambda=e^{-n^2/(2\lambda)}.
\]

Since \(|d_m(t)-1|\leq1/4\),

\[
\operatorname{Re}\frac1{\lambda_m(t)}
 =\frac1{\alpha\kappa_m}\operatorname{Re}\frac1{d_m(t)}
 \geq\frac1{3\alpha\kappa_m}.
\]

Normal convergence of the Fourier series permits termwise integration.
Equations (22) and the preceding estimate imply

\[
\tag{23}
 \left|
 \langle G_m(t,\cdot)\rangle_{\lambda_m(t)}
       -\left\langle G_m(t,\cdot)\right\rangle_{\mathbb T}\right|
 \leq K_\alpha e^{-c_\alpha/\kappa_m}
 \qquad(|t|\leq1/2),
\]

where one may take \(c_\alpha=1/(6\alpha)\). The same estimate, with
another constant, holds after multiplying by \(d_m(t)^\nu/(\alpha-1)\).
Cauchy's estimate bounds the resulting \(m\)th derivative at zero by

\[
K_\alpha 2^m m! e^{-c_\alpha/\kappa_m}=o_\alpha(\beta^m),
\]

because \(1/\kappa_m=2^{m/4}/8\). This construction also supplies the
stated holomorphic extension, agreeing with (20) on real nonnegative
times.

\paragraph{The constant periodic
mode.}\label{the-constant-periodic-mode.}

The zero Fourier coefficient is

\[
\left\langle G_m(t,\cdot)\right\rangle_{\mathbb T}=1+(\alpha-1)D_{\alpha,m}(\tau_m(t)).
\]

Thus, after division by \(\alpha-1\), the constant term contributes
\(d_m(t)^\nu/(\alpha-1)\). Its \(m\)th derivative has absolute value at
most

\[
\frac{(|\nu|+m)^m\kappa_m^m}{|\alpha-1|}
 =\exp\{-\tfrac14(\log2)m^2+O_\alpha(m\log m)\}
 =o_\alpha(\beta^m).
\]

\paragraph{Time change and the remaining scalar
factor.}\label{time-change-and-the-remaining-scalar-factor.}

We use a uniform estimate for each spectral mode. For \(L\geq1\) and
integer \(n\geq1\), define

\[
T_n(L,\kappa)=
 (-1)^n\left.\frac{\,\mathrm d^n}{\,\mathrm dt^n}
        e^{-Lt/(1+\kappa t)}\right|_{t=0}.
\]

Expanding the exponential and then \((1+\kappa t)^{-\ell}\) yields

\[
\tag{24}
 T_n(L,\kappa)=
 \sum_{r=0}^{n-1}\frac{n!}{(n-r)!}
       \binom{n-1}{r}\kappa^r L^{n-r},
 \qquad T_0(L,\kappa)=1.
\]

This identity and the associated estimate are the time-change
calculation of~\href{https://arxiv.org/abs/2608.30275v1}{{[}3{]}}. Every
term is nonnegative, and

\[
\tag{25}
 1\leq T_n(L,\kappa)/L^n\leq e^{n^2\kappa/L}
 \quad(n\geq0),
\]

where the lower bound for \(n=0\) is an equality.

For the extra factor required by the power moment, put

\[
U_m(L,\kappa)=(-1)^m\left.\frac{\,\mathrm d^m}{\,\mathrm dt^m}
  \left[(1+\kappa t)^\nu e^{-Lt/(1+\kappa t)}\right]\right|_{t=0}.
\]

Leibniz' rule gives

\[
U_m=
 \sum_{k=0}^m\binom mk(-1)^k(\nu )_{\underline{k}}\kappa^k
       T_{m-k}(L,\kappa).
\]

Using \(|(\nu )_{\underline{k}}|\leq(|\nu|+m)^k\),
\(\binom mk\leq m^k/k!\), and (25), we obtain

\[
\tag{26}
 \left|\frac{U_m(L,\kappa)}{L^m}-1\right|
 \leq
 \exp\!\left(\frac{(2m^2+|\nu|m)\kappa}{L}\right)-1.
\]

For clarity, the \(k=0\) term contributes at most \(e^{m^2\kappa/L}-1\)
to this difference; the sum over \(k\geq1\) is at most

\[
e^{m^2\kappa/L}
   \big(e^{m(|\nu|+m)\kappa/L}-1\big),
\]

whose sum is (26).

For \(\kappa=\kappa_m\), the right side of (26) is \(o(1)\) uniformly in
\(L\geq1\). The periodic expansion can be written

\[
D_{\alpha,m}(s)=\sum_{j\geq1}A_j(\alpha)2^{-mj/2}e^{-js}.
\]

Apply (26) with \(L=j\) and use the absolute-sum estimate in Lemma~3. It
follows that

\[
(-1)^m\left.\frac{\,\mathrm d^m}{\,\mathrm dt^m}
       [d_m(t)^\nu D_{\alpha,m}(\tau_m(t))]\right|_{t=0}
 =(-1)^mD_{\alpha,m}^{(m)}(0)+o_\alpha(\beta^m).
\]

Termwise differentiation is justified by the same absolutely convergent
monomial series as in Lemma~3, or by (15) and (26). Combining the three
estimates proves (21).~\(\square\)

\emph{Proof of Theorem~2.} Positivity, Gaussian tails, smoothness and
strict log-concavity were proved in (16). From (1),

\[
(-1)^{m-1}S_{\alpha,f}^{(m)}(t)
 =\frac{(-1)^mJ_{\alpha,f}^{(m)}(t)}{\alpha-1}.
\]

Combine Lemmas~3 and~4 to obtain (8). Its leading coefficient is
strictly negative. The holomorphic extension in Lemma~4 gives continuity
of the \(m\)th derivative at zero, so the strict sign persists for a
sufficiently small interval of positive times.~\(\square\)

\section{The quadratic case, Shannon order and higher
dimensions}\label{the-quadratic-case-shannon-order-and-higher-dimensions}

\textbf{Proposition 5} (The quadratic case). For any probability density
\(f\) on \(\mathbb R^d\), \(t>0\) and integer \(k\geq1\),

\[
\tag{27}
 (-1)^{k-1}S_{2,f}^{(k)}(t)
 =\frac1{(2\pi)^d}\int_{\mathbb R^d}|\xi|^{2k}
    |\widehat f(\xi)|^2e^{-t|\xi|^2}\,\mathrm d\xi\geq0.
\]

In dimension one this equals
\(\int_\mathbb R|\partial_x^kP_tf|^2\,\mathrm dx\).

\emph{Proof.} Use the Fourier convention
\(\widehat f(\xi)=\int e^{-ix\cdot\xi}f(x)\,\mathrm dx\). Since
\(|\widehat f|\leq1\) and \(\widehat{P_tf}=e^{-t|\xi|^2/2}\widehat f\),
Plancherel's identity gives

\[
J_{2,f}(t)=\frac1{(2\pi)^d}\int|\widehat f(\xi)|^2e^{-t|\xi|^2}\,\mathrm d\xi.
\]

Gaussian domination permits differentiation of every order at positive
times. Since \(S_{2,f}=1-J_{2,f}\), (27) follows. This reproves the
quadratic fact in~\href{https://arxiv.org/abs/2312.01819v4}{{[}1{]}} and
its arbitrary-dimensional version
in~\href{https://arxiv.org/abs/2504.01837v2}{{[}4, Appendix
B{]}}.~\(\square\)

At the Shannon order, we use the already established theorem of Zou et
al.~\href{https://arxiv.org/abs/2608.30275v1}{{[}3, Theorem 1{]}}. With
\(a_m=2^{-m/4}\), \(b_m=2^{-m}\) and \(\kappa_m=8a_m\), their density
(5) is smooth, strictly positive and strictly log-concave for
\(m\geq5\), and

\[
\tag{28}
 (-1)^{m-1}S_{1,f_m}^{(m)}(0)
 =-\frac{11}{96}\beta^m+o(\beta^m).
\]

Their analytic construction also gives failure at small positive times.
Equation (28) is quoted from prior work, not asserted as a new case of
our result. Together with Theorem~2 and Proposition~5, it proves
Theorem~1 in dimension one.

\textbf{Corollary 6} (Every dimension). For every \(d\geq2\) and
\(\alpha>0\), \(\alpha\ne2\), there is a smooth strictly positive
strictly log-concave probability density with Gaussian tails on
\(\mathbb R^d\) for which (2) fails.

\emph{Proof.} First take a one-dimensional density \(f\), a derivative
order \(k\), and a positive time \(t_0\) with a strict negative signed
derivative, as obtained above. Let \(g_R\) be the centred Gaussian
density of covariance \(RI_{d-1}\) and use \(f\otimes g_R\). Strict
log-concavity and the other stated regularity properties are preserved.

For \(\alpha\ne1\), the power moments factor, and

\[
J_{\alpha,g_R}(t)=K_{\alpha,R,d}
  (1+t/R)^{-(d-1)(\alpha-1)/2},
\]

where \(K_{\alpha,R,d}>0\) is independent of time. The factor following
\(K_{\alpha,R,d}\) tends to \(1\) in \(C^k\) on every compact time
interval as \(R\to\infty\). Consequently the signed \(k\)th derivative
of \(S_{\alpha,f\otimes g_R}\), after division by \(K_{\alpha,R,d}\),
converges at \(t_0\) to the strictly negative signed derivative for
\(f\). It is therefore negative for sufficiently large \(R\).

For \(\alpha=1\), Shannon entropy is additive, and the Gaussian summand
contributes

\[
(-1)^{k-1}S_{1,g_R}^{(k)}(t_0)
 =\frac{d-1}{2}\frac{(k-1)!}{(R+t_0)^k}\longrightarrow0.
\]

Thus the strict negative sign is again preserved for large \(R\).
Combining this corollary with Proposition~5 completes Theorem~1 in every
dimension.~\(\square\)

\section{\texorpdfstring{A completely specified witness at order
\(\alpha=3/2\)}{A completely specified witness at order \textbackslash alpha=3/2}}\label{a-completely-specified-witness-at-order-alpha32}

The asymptotic family proves the classification and the log-concave
strengthening, but it does not give a numerical threshold for \(m\) at
each \(\alpha\). The following separate construction gives a finite
certificate at derivative order \(13\) for the open-interval order
\(\alpha=3/2\).

\textbf{Theorem 7} (Finite real-line witness). Let \(R=2^{40}\) and let
\(\gamma_R\) denote the centred Gaussian density of variance \(R\). Set

\[
\tag{29}
 \begin{aligned}
 f_R(x)&=\frac{\gamma_R(x)}{Z_R}
    \left(1+\frac1{10}\cos x+\frac1{32000}\cos2x\right),\\
 Z_R&=1+\frac1{10}e^{-R/2}+\frac1{32000}e^{-2R}.
 \end{aligned}
\]

This is a smooth strictly positive probability density with Gaussian
tails. Under the convention \(P_t=e^{t\Delta/2}\), put

\[
C_R=(3/2)^{-1/2}(2\pi R)^{-1/4}.
\]

Then

\[
\tag{30}
 S_{3/2,f_R}^{(13)}(1/50)<-\frac{C_R}{1024}<0.
\]

Since \((-1)^{12}=1\), this violates (2). No log-concavity assertion is
made for this particular finite witness.

\emph{Proof.} For the finite calculation only, write
\(Q_t=e^{t\Delta}=P_{2t}\). On the circle define

\[
v(t,x)=1+ae^{-t}\cos x+be^{-4t}\cos2x,
 \qquad a=\frac1{10},\quad b=\frac1{32000},\quad
 F(t)=\left\langle v(t,\cdot)^{3/2}\right\rangle_{\mathbb T}.
\]

The exact rational certificate in Appendix~A proves

\[
\tag{31}
 (-1)^{13}F^{(13)}(1/100)<-6.
\]

We give all analytic estimates needed to transfer this finite sign to
(29).

Put \(M_R(t)=\int(Q_tf_R)^{3/2}\,\mathrm dx\), \(d=1+2t/R\) and
\(\tau=t/d\). The exact window identity gives

\[
Q_tf_R(x)=Z_R^{-1}\gamma_{Rd}(x)
  [1+ae^{-\tau}\cos(x/d)+be^{-4\tau}\cos(2x/d)].
\]

Let \(\widehat V_\ell(s)\) be the Fourier coefficients in \(x\) of
\(v(s,x)^{3/2}\). Gaussian integration gives

\[
\tag{32}
 G_R(t):=\frac{M_R(t)}{C_R}
 =Z_R^{-3/2}d^{-1/4}
   \sum_{\ell\in\mathbb Z}\widehat V_\ell(\tau)
                   e^{-\ell^2R/(3d)}.
\]

In particular, \(\widehat V_0=F\). The normally convergent Fourier
expression on the right, rather than a complex power of the real-line
density, supplies the holomorphic extension used below.

Let \(T=1/100\) and use the disk \(|z-T|\leq1\). The following bounds
hold for every \(R\geq1024\), hence for the specified \(R\). We have
\(|z|<11/10\),

\[
|d-1|\leq\frac{11}{5R},\qquad
 \operatorname{Re}(1/d)>\frac12,\qquad
 |\tau-z|\leq\frac{121}{50R(1-11/(5R))}<\frac5{2R}.
\]

Every point on the segment from \(z\) to \(\tau\) has real part greater
than \(-11/10\). For \(\operatorname{Re}s\geq-11/10\) and
\(|\operatorname{Im}x|\leq1\), use

\[
e^{11/10}<10/3,\qquad e^{22/5}<135,\qquad
 \cosh1<2,\qquad \cosh2<4.
\]

For example, the exponential bounds follow from \(e<3\),
\(e^{1/10}<10/9\) and \(e^{2/5}<5/3\), the last two being immediate
geometric majorants of the exponential series. Hence

\[
|ae^{-s}\cos x+be^{-4s}\cos2x|
 <\frac23+\frac{27}{1600}<\frac34.
\]

The principal \(3/2\) power is therefore holomorphic on this domain, has
modulus less than \(3\), and obeys

\[
|\widehat V_\ell(s)|\leq3e^{-|\ell|}.
\]

Also \(|F|<3\), and differentiation under the circle average gives

\[
|F'(s)|\leq\frac32\sqrt{\frac74}
       \left(\frac23+\frac{27}{400}\right)<2.
\]

Consequently \(|F(\tau)-F(z)|<5/R\). The nonzero Fourier modes in (32)
satisfy

\[
\tag{33}
 \left|\sum_{\ell\ne0}\widehat V_\ell(\tau)
                  e^{-\ell^2R/(3d)}\right|
 <6e^{-R/6}.
\]

Integration of the derivative of \(w^{-1/4}\) along the segment from
\(1\) to \(d\) gives

\[
|d^{-1/4}-1|<\frac{11}{10R},\qquad |d^{-1/4}|<2.
\]

Furthermore, \(Z_R\geq1\) and

\[
Z_R-1<\frac{11}{100}e^{-R/2}\leq\frac{11}{100R}.
\]

Since \(|Z_R^{-3/2}-1|\leq(3/2)(Z_R-1)\), the factor
\(A=Z_R^{-3/2}d^{-1/4}\) obeys

\[
|A-1|<\frac2R,\qquad |A|<2.
\]

For \(R\geq1024\), \(12e^{-R/6}<1/R\): the positive exponential series
gives \(e^{R/6}>R^3/1296\), and \(R^2>15552\). Writing

\[
A F(\tau)-F(z)=[F(\tau)-F(z)]+(A-1)F(\tau)
\]

and using (33), the time-change, scalar-factor and nonzero-mode errors
are respectively smaller than \(5/R\), \(6/R\) and \(1/R\). We may
therefore use the generous bound

\[
\tag{34}
 \sup_{|z-T|\leq1}|G_R(z)-F(z)|<100/R.
\]

The Fourier series in (32) converges normally on this disk, so it
defines a holomorphic function agreeing with \(M_R/C_R\) at real
nonnegative times. Cauchy's estimate gives

\[
\tag{35}
 |G_R^{(13)}(T)-F^{(13)}(T)|
 <\frac{100\,13!}{2^{40}}
 =\frac{152026875}{268435456}<1.
\]

The final comparison is also checked in exact integer arithmetic by the
accompanying certificate.

Equation (31) says \(F^{(13)}(T)>6\), so (35) implies
\(M_R^{(13)}(T)/C_R>5\). Finally, \(P_t=Q_{t/2}\), \(t=2T=1/50\), and
\(S_{3/2,f}=2(1-J_{3/2,f})\). Thus

\[
J_{3/2,f_R}^{(13)}(1/50)>\frac{5C_R}{8192},
 \qquad
 S_{3/2,f_R}^{(13)}(1/50)<-\frac{5C_R}{4096}
                         <-\frac{C_R}{1024}.
\]

This proves the claim.~\(\square\)

\section{Dilation and the scale of a
counterexample}\label{dilation-and-the-scale-of-a-counterexample}

\textbf{Corollary 8} (Elementary dilation consequences). Let \(f\) be a
one-dimensional density with finite positive variance, and suppose its
signed entropy derivative of order \(k\geq1\) is strictly negative at a
positive time \(t_0\). For \(a>0\) define \(f_a(x)=a^{-1}f(x/a)\). Then
the same signed derivative for \(f_a\) is strictly negative at
\(a^2t_0\). Consequently a witness can be rescaled to variance one, or
it can be rescaled to have failure at any prescribed positive time.
These are separate choices of the single scale parameter.

\emph{Proof.} Gaussian convolution and the substitution \(x=ay\) give
\[
P_t f_a(x)=a^{-1}(P_{t/a^2}f)(x/a).
\]
For \(\alpha\ne1\) this implies
\[
J_{\alpha,f_a}(t)=a^{1-\alpha}J_{\alpha,f}(t/a^2),
\]
\[
(-1)^{k-1}S_{\alpha,f_a}^{(k)}(t)
=a^{1-\alpha-2k}(-1)^{k-1}S_{\alpha,f}^{(k)}(t/a^2).
\]
The multiplicative factor is positive. For \(\alpha=1\),
\(S_{1,f_a}(t)=S_{1,f}(t/a^2)+\log a\), so the derivative factor is
\(a^{-2k}\), again positive. Variance transforms as
\(\operatorname{Var}(f_a)=a^2\operatorname{Var}(f)\). Thus
\(a=\operatorname{Var}(f)^{-1/2}\) gives variance one, whereas
\(a=(T/t_0)^{1/2}\) places failure at a prescribed \(T>0\).

Smoothness, strict positivity and Gaussian tails are preserved by
dilation. If \(f\) is strictly log-concave, then
\(\log f_a(x)=\log f(x/a)-\log a\) is strictly concave by affine composition.
For the constructed family its negative curvature transforms by
\((\log f_a)''(x)=a^{-2}(\log f)''(x/a)\). The constructed densities have finite positive variance.
This proves the conclusions. \(\square\)

This corollary is a standard heat-scaling consequence, not a new
localisation mechanism. The ratio of failure time to variance is
unchanged: the two quantities cannot be prescribed independently.
Dilation does not supply a log-concavity claim for the finite witness
and does not establish a uniform curvature bound across the asymptotic
family.

\section{What the classification
establishes}\label{what-the-classification-establishes}

Theorem~1 settles the universal order question, including each
individual order in \((1,2)\). The exception at \(\alpha=2\) reflects
the quadratic Fourier representation (27). For \(\alpha\notin\{1,2\}\),
the power-entropy expansion has a nonzero cubic contribution; the second
harmonic can be chosen to make the combined dominant spectral
coefficient negative. The Shannon case is supplied separately by its
cited entropy expansion. The sixth-order contribution at the same
spectral level is retained exactly in (14), rather than discarded as a
lower-order perturbation.

The bounds are pointwise in \(\alpha\). In the displayed family, the
sufficient curvature lower bound \(5a_m\) tends to zero; no uniform
curvature or variance bound across that family is claimed. Dilation
gives a different normalization of each witness, not a freely chosen
physical noise scale.

The result does not determine the first derivative order at which a
particular \(\alpha\) admits failure. The asymptotic family establishes
existence with strong regularity and shape restrictions, while the
order-\(13\) witness gives one reproducible finite example. Neither
construction estimates a physically typical noise level. The method and
its limitations are inherited in substantial part from the Shannon
construction of~\href{https://arxiv.org/abs/2608.30275v1}{{[}3{]}}; the
contribution here is the complete positive-order classification and the
accompanying power-moment and finite-certificate details.

\section*{Appendix A. Exact rational circle
certificate}\label{exact-rational-circle-certificate}

This appendix explains the finite certificate used in (31). It involves
no floating-point sign decision. Let

\[
\mathcal F(z)=\left\langle (1+az\cos x+bz^4\cos 2x)^{3/2}\right\rangle_{\mathbb T}
             =\sum_{k\geq0}c_kz^k,
 \qquad a=1/10,\quad b=1/32000.
\]

Then \(F(t)=\mathcal F(e^{-t})\). The coefficients are rational and can
be computed as

\[
\tag{36}
 c_k=\sum_{i+4j=k}
 \binom{3/2}{i+j}\binom{i+j}{i}a^ib^j
      \left\langle \cos^i x\cos^j2x\right\rangle_{\mathbb T}.
\]

The moments in this formula are exactly

\[
\left\langle \cos^i x\cos^j2x\right\rangle_{\mathbb T}
 =2^{-(i+j)}
  \sum_{\substack{0\leq r\leq i,\,0\leq s\leq j\\
           i-2r+2(j-2s)=0}}\binom ir\binom js.
\]

Odd coefficients vanish by the substitution \(x\mapsto\pi-x\). The
script computes every coefficient through degree \(23\); the first four
nonzero coefficients after \(c_0=1\) are

\[
c_2=\frac3{1600},\quad
 c_4=\frac9{10240000},\quad
 c_6=-\frac{41}{3276800000},\quad
 c_8=\frac{12213}{83886080000000}.
\]

At derivative order \(13\) and time \(T=1/100\),

\[
\tag{37}
 (-1)^{13}F^{(13)}(T)=\sum_{k\geq1}c_k k^{13}e^{-kT}.
\]

For the infinite tail, use the analytic radius \(\rho=8\). Since

\[
a\rho+b\rho^4=\frac{116}{125}<1,
 \qquad (1+116/125)^3<3^2,
\]

the principal power is analytic for \(|z|\leq8\) and
\(|\mathcal F(z)|<3\) there. Thus \(|c_k|\leq3\,8^{-k}\). The ratio of
successive even Cauchy majorants for \(k\geq24\) is at most

\[
q=\frac{(13/12)^{13}}{64}<1.
\]

Since \(e^{-kT}\leq1\), the absolute tail in (37) is bounded by

\[
\tag{38}
 H=\frac{3\,24^{13}}{8^{24}(1-q)}<\frac1{1000}.
\]

To bound the finite sum with rational arithmetic, bracket
\(E=e^{1/100}\) by

\[
E_-=\sum_{j=0}^{12}\frac{(1/100)^j}{j!},
 \qquad
 E_+=E_-+\frac{(1/100)^{13}/13!}{1-1/1400}.
\]

The omitted positive Taylor terms have successive ratio at most
\(1/1400\). Therefore the finite sum in (37) is at most

\[
B=
 \sum_{\substack{2\leq k\leq22\\ k\ \mathrm{even}\\c_k\geq0}}
       \frac{c_k k^{13}}{E_-^k}
 +
 \sum_{\substack{2\leq k\leq22\\ k\ \mathrm{even}\\c_k<0}}
       \frac{c_k k^{13}}{E_+^k}.
\]

The script checks the exact rational inequalities

\[
B<-\frac{27}{4},\qquad H<\frac1{1000}.
\]

Consequently

\[
(-1)^{13}F^{(13)}(1/100)\leq B+H
 <-\frac{6749}{1000}<-6,
\]

which proves (31). The readable diagnostic for \(B+H\) is approximately
\(-6.753909793468304\); it is not used to decide the sign.

The file \texttt{checks/tsallis\_order13\_certificate.py} contains the
exact certificate. It computes the coefficients, verifies the
analytic-bound hypotheses and finite rational inequalities, and checks
the integer comparison in (35) and the final entropy margin.
The companion \texttt{checks/test\_guards.py} checks intended failures.
Each accepts \texttt{--output NEW\_PATH.json} and refuses to overwrite
an existing receipt.
The mathematical certificate above is complete without the source
listing. Executable acceptance uses explicit named guards rather than
removable assertions; see \texttt{VERIFICATION.md} for the exact
invocation, positive and intended-failure checks, interpreter modes and
artifact identities. Arithmetic success verifies only the finite
comparisons, not the analytic convolution, Fourier or Cauchy arguments.

\section*{Declarations}\label{declarations}

\textbf{Data and code availability.} This theoretical study uses no
empirical dataset. The complete mathematical certificate is included
above; the accompanying public package contains its executable checker
and verification documentation.

\textbf{AI assistance and research status.} AI systems assisted the
analysis, manuscript revision and internal checks. This is an unrefereed
candidate, not external human peer review, unaffiliated reproduction or
proof-assistant verification.

\textbf{Contributions.} Anonymous authorship is retained. The work
comprises mathematical analysis, software, validation and writing; no
individual human CRediT allocation is asserted from the available
record.

\textbf{Funding and competing interests.} No funding or
competing-interest information was supplied with the research record;
absence of a declaration is not a verified assertion that none exists.

\textbf{Ethics.} No human participants, animal research or personal data
are involved in the mathematical material presented.

\section*{References}\label{references}

\paragraph{{[}1{]}}\label{section}

H. Wu, L. Yu and L. Guo, \emph{On the Complete Monotonicity of Rényi
Entropy}, IEEE Transactions on Information Theory \textbf{71}(7) (2025),
4895--4914.
\href{https://doi.org/10.1109/TIT.2025.3561384}{doi:10.1109/TIT.2025.3561384}.
\href{https://arxiv.org/abs/2312.01819v4}{arXiv:2312.01819}. The Tsallis
conjecture is Conjecture 1, Section II-C, equation (19) in arXiv version
4, dated 1 November 2025.

\paragraph{{[}2{]}}\label{section-1}

Y. Gu and M. Sellke, \emph{A Counterexample to the Gaussian Completely
Monotone Conjecture}, 2026 preprint.
\href{https://arxiv.org/abs/2605.11656v1}{arXiv:2605.11656}.

\paragraph{{[}3{]}}\label{section-2}

J. Zou, L. Fan, J. Gao and J. Wang, \emph{An Explicit Family of
Log-Concave Counterexamples to the Gaussian Completely Monotone
Conjecture}, 2026 preprint, version 1, 31 August 2026.
\href{https://arxiv.org/abs/2608.30275v1}{arXiv:2608.30275}.

\paragraph{{[}4{]}}\label{section-3}

H. Wu and L. Yu, \emph{Entropic Isoperimetric and Cramér--Rao
Inequalities for Rényi--Fisher Information}, 2025 preprint, version 2,
11 June 2025.
\href{https://arxiv.org/abs/2504.01837v2}{arXiv:2504.01837}.

\paragraph{{[}5{]}}\label{section-4}

L. Sun, \emph{Concavity of Tsallis Entropy and Tsallis Entropy Power
along Heat Flow}, 2026 preprint, version 3, 12 June 2026.
\href{https://arxiv.org/abs/2601.12944v3}{arXiv:2601.12944}.

\end{document}
